\section{Dataset Statistics}
\label{appx:dataset}
We present additional breakdowns and analyses of the \bugs{} dataset, focusing on the kinds of repositories and bugs that are represented.

\textbf{Repository categorization.}
We present an exhaustive list of repositories used in \bugs{} in Table~\ref{tab:summary_of_repos}.
We categorize the repositories into seven general buckets: Data Parsing and Transformation ($39$), Web \& API Development ($11$), Code Quality \& Testing ($12$), Visualization \& Presentation ($8$), System Tools \& Protocols ($17$), Natural Language Processing ($7$), and Miscellaneous ($6$).
The categorizations were performed by first, determining an appropriate set of categories based on manual inspection supported by the descriptions and GitHub topics associated with each repository.
After settling upon the buckets, we asked GPT-4o to provide a label based on the repository's metadata and \texttt{README} dump.
\bugs{} represents a wider and more variegated coverage of software tools and applications compared to any prior works.

\begin{CJK*}{UTF8}{gbsn}
\begin{longtable}{p{0.3\linewidth} | p{0.65\linewidth}}
\toprule
\textbf{Repository} & \textbf{Description} \\
\midrule
\endfirsthead

\toprule
\textbf{Repository} & \textbf{Description} \\
\midrule
\endhead

\midrule
\multicolumn{2}{r}{\textit{Continued on next page}} \\
\endfoot

\bottomrule
\endlastfoot

\multicolumn{2}{c}{\textit{Code Quality and Testing}} \\
\midrule
PyCQA/flake8 & flake8 is a python tool that glues together pycodestyle, pyflakes, mccabe, and third-party plugins to check the style and quality of some python code. \\
Suor/funcy & A fancy and practical functional tools \\
adrienverge/yamllint & A linter for YAML files. \\
agronholm/typeguard & Run-time type checker for Python \\
cknd/stackprinter & Debugging-friendly exceptions for Python \\
cool-RR/PySnooper & Never use print for debugging again \\
getmoto/moto & A library that allows you to easily mock out tests based on AWS infrastructure. \\
pylint-dev/astroid & A common base representation of python source code for pylint and other projects \\
pytest-dev/iniconfig None
pytest-dev/iniconfig & None \\
python/mypy & Optional static typing for Python \\
pyupio/safety & Safety checks Python dependencies for known security vulnerabilities and suggests the proper remediations for vulnerabilities detected. \\
pyutils/line\_profiler & Line-by-line profiling for Python \\
rubik/radon & Various code metrics for Python code \\
spulec/freezegun & Let your Python tests travel through time \\
sqlfluff/sqlfluff & A modular SQL linter and auto-formatter with support for multiple dialects and templated code. \\
\midrule
\multicolumn{2}{c}{\textit{Data Parsing and Transformation}} \\
\midrule
alecthomas/voluptuous & CONTRIBUTIONS ONLY: Voluptuous, despite the name, is a Python data validation library. \\
andialbrecht/sqlparse & A non-validating SQL parser module for Python \\
buriy/python-readability & fast python port of arc90's readability tool, updated to match latest readability.js! \\
burnash/gspread & Google Sheets Python API \\
chardet/chardet & Python character encoding detector \\
cloudpipe/cloudpickle & Extended pickling support for Python objects \\
dask/dask & Parallel computing with task scheduling \\
datamade/usaddress & :us: a python library for parsing unstructured United States address strings into address components \\
davidhalter/parso & A Python Parser \\
erikrose/parsimonious & The fastest pure-Python PEG parser I can muster \\
facelessuser/soupsieve & A modern CSS selector implementation for BeautifulSoup \\
gawel/pyquery & A jquery-like library for python \\
google/textfsm & Python module for parsing semi-structured text into python tables. \\
gruns/furl & �� URL parsing and manipulation made easy. \\
gweis/isodate & ISO 8601 date/time parser \\
hukkin/tomli & A lil' TOML parser \\
jawah/charset\_normalizer & Truly universal encoding detector in pure Python \\
john-kurkowski/tldextract & Accurately separates a URL’s subdomain, domain, and public suffix, using the Public Suffix List (PSL). \\
joke2k/faker & Faker is a Python package that generates fake data for you. \\
jsvine/pdfplumber & Plumb a PDF for detailed information about each char, rectangle, line, et cetera — and easily extract text and tables. \\
kayak/pypika & PyPika is a python SQL query builder that exposes the full richness of the SQL language using a syntax that reflects the resulting query. PyPika excels at all sorts of SQL queries but is especially useful for data analysis. \\
keleshev/schema & Schema validation just got Pythonic \\
kennethreitz/records & SQL for Humans™ \\
kurtmckee/feedparser & Parse feeds in Python \\
lepture/mistune & A fast yet powerful Python Markdown parser with renderers and plugins. \\
madzak/python-json-logger & Json Formatter for the standard python logger \\
mahmoud/glom & ☄️ Python's nested data operator (and CLI), for all your declarative restructuring needs. Got data? Glom it! ☄️ \\
marshmallow-code/marshmallow & A lightweight library for converting complex objects to and from simple Python datatypes. \\
martinblech/xmltodict & Python module that makes working with XML feel like you are working with JSON \\
matthewwithanm/python-markdownify & Convert HTML to Markdown \\
mewwts/addict & The Python Dict that's better than heroin. \\
mido/mido & MIDI Objects for Python \\
modin-project/modin & Modin: Scale your Pandas workflows by changing a single line of code \\
mozilla/bleach & Bleach is an allowed-list-based HTML sanitizing library that escapes or strips markup and attributes \\
msiemens/tinydb & TinyDB is a lightweight document oriented database optimized for your happiness :) \\
pandas-dev/pandas & Flexible and powerful data analysis / manipulation library for Python, providing labeled data structures similar to R data.frame objects, statistical functions, and much more \\
pdfminer/pdfminer.six & Community maintained fork of pdfminer - we fathom PDF \\
pudo/dataset & Easy-to-use data handling for SQL data stores with support for implicit table creation, bulk loading, and transactions. \\
pydantic/pydantic & Data validation using Python type hints \\
pydata/patsy & Describing statistical models in Python using symbolic formulas \\
pydicom/pydicom & Read, modify and write DICOM files with python code \\
pygments/pygments & Pygments is a generic syntax highlighter written in Python \\
pyparsing/pyparsing & Python library for creating PEG parsers \\
python-jsonschema/jsonschema & An implementation of the JSON Schema specification for Python \\
python-openxml/python-docx & Create and modify Word documents with Python \\
r1chardj0n3s/parse & Parse strings using a specification based on the Python format() syntax. \\
scanny/python-pptx & Create Open XML PowerPoint documents in Python \\
scrapy/scrapy & Scrapy, a fast high-level web crawling \& scraping framework for Python. \\
seperman/deepdiff & DeepDiff: Deep Difference and search of any Python object/data. DeepHash: Hash of any object based on its contents. Delta: Use deltas to reconstruct objects by adding deltas together. \\
sloria/environs & simplified environment variable parsing \\
sunpy/sunpy & SunPy - Python for Solar Physics \\
tkrajina/gpxpy & gpx-py is a python GPX parser. GPX (GPS eXchange Format) is an XML based file format for GPS tracks. \\
tobymao/sqlglot & Python SQL Parser and Transpiler \\
un33k/python-slugify & Returns unicode slugs \\
\midrule
\multicolumn{2}{c}{\textit{Machine Learning and AI}} \\
\midrule
facebookresearch/fvcore & Collection of common code that's shared among different research projects in FAIR computer vision team. \\
facebookresearch/hydra & Hydra is a framework for elegantly configuring complex applications \\
HIPS/autograd & Efficiently computes derivatives of NumPy code. \\
iterative/dvc & �� Data Versioning and ML Experiments \\
jaraco/inflect & Correctly generate plurals, ordinals, indefinite articles; convert numbers to words \\
life4/textdistance & �� Compute distance between sequences. 30+ algorithms, pure python implementation, common interface, optional external libs usage. \\
luozhouyang/python-string-similarity & A library implementing different string similarity and distance measures using Python. \\
Mimino666/langdetect & Port of Google's language-detection library to Python. \\
mozillazg/python-pinyin & 汉字转拼音(pypinyin) \\
pndurette/gTTS & Python library and CLI tool to interface with Google Translate's text-to-speech API \\
Project-MONAI/MONAI & AI Toolkit for Healthcare Imaging \\
seatgeek/thefuzz & Fuzzy String Matching in Python \\
vi3k6i5/flashtext & Extract Keywords from sentence or Replace keywords in sentences. \\
\midrule
\multicolumn{2}{c}{\textit{System Tools and Protocols}} \\
\midrule
agronholm/exceptiongroup & Backport of PEP 654 (exception groups) \\
aio-libs/async-timeout & asyncio-compatible timeout class \\
arrow-py/arrow & �� Better dates \& times for Python \\
borntyping/python-colorlog & A colored formatter for the python logging module \\
cantools/cantools & CAN bus tools. \\
conan-io/conan & Conan - The open-source C and C++ package manager \\
cookiecutter/cookiecutter & A cross-platform command-line utility that creates projects from cookiecutters (project templates), e.g. Python package projects, C projects. \\
dbader/schedule & Python job scheduling for humans. \\
gruns/icecream & �� Never use print() to debug again. \\
jd/tenacity & Retrying library for Python \\
mahmoud/boltons & �� Like builtins, but boltons. 250+ constructs, recipes, and snippets which extend (and rely on nothing but) the Python standard library.  Nothing like Michael Bolton. \\
oauthlib/oauthlib & A generic, spec-compliant, thorough implementation of the OAuth request-signing logic \\
pallets/click & Python composable command line interface toolkit \\
paramiko/paramiko & The leading native Python SSHv2 protocol library. \\
pexpect/ptyprocess & Run a subprocess in a pseudo terminal \\
pyasn1/pyasn1 & Generic ASN.1 library for Python \\
pyca/pyopenssl & A Python wrapper around the OpenSSL library \\
python-hyper/h11 & A pure-Python, bring-your-own-I/O implementation of HTTP/1.1 \\
python-trio/trio & Trio – a friendly Python library for async concurrency and I/O \\
rustedpy/result & NOT MAINTAINED - A simple Rust like Result type for Python 3. Fully type annotated. \\
termcolor/termcolor & ANSI color formatting for output in terminal \\
theskumar/python-dotenv & Reads key-value pairs from a .env file and can set them as environment variables. It helps in developing applications following the 12-factor principles. \\
tox-dev/pipdeptree & A command line utility to display dependency tree of the installed Python packages \\
\midrule
\multicolumn{2}{c}{\textit{Visualization and Presentation}} \\
\midrule
amueller/word\_cloud & A little word cloud generator in Python \\
lincolnloop/python-qrcode & Python QR Code image generator \\
prettytable/prettytable & Display tabular data in a visually appealing ASCII table format \\
pwaller/pyfiglet & An implementation of figlet written in Python \\
rsalmei/alive-progress & A new kind of Progress Bar, with real-time throughput, ETA, and very cool animations! \\
weaveworks/grafanalib & Python library for building Grafana dashboards \\
\midrule
\multicolumn{2}{c}{\textit{Web and API Development}} \\
\midrule
Cog-Creators/Red-DiscordBot & A multi-function Discord bot \\
Knio/dominate & Dominate is a Python library for creating and manipulating HTML documents using an elegant DOM API.  It allows you to write HTML pages in pure Python very concisely, which eliminate the need to learn another template language, and to take advantage of the more powerful features of Python. \\
alanjds/drf-nested-routers & Nested Routers for Django Rest Framework \\
benoitc/gunicorn & gunicorn 'Green Unicorn' is a WSGI HTTP Server for UNIX, fast clients and sleepy applications. \\
bottlepy/bottle & bottle.py is a fast and simple micro-framework for python web-applications. \\
django-money/django-money & Money fields for Django forms and models. \\
django/channels & Developer-friendly asynchrony for Django \\
django/daphne & Django Channels HTTP/WebSocket server \\
encode/starlette & The little ASGI framework that shines. �� \\
getnikola/nikola & A static website and blog generator \\
graphql-python/graphene & GraphQL framework for Python \\
marshmallow-code/apispec & A pluggable API specification generator. Currently supports the OpenAPI Specification (f.k.a. the Swagger specification).. \\
marshmallow-code/webargs & A friendly library for parsing HTTP request arguments, with built-in support for popular web frameworks, including Flask, Django, Bottle, Tornado, Pyramid, webapp2, Falcon, and aiohttp. \\
pallets/jinja & A very fast and expressive template engine. \\
pallets/markupsafe & Safely add untrusted strings to HTML/XML markup. \\
tornadoweb/tornado & Tornado is a Python web framework and asynchronous networking library, originally developed at FriendFeed. \\
tweepy/tweepy & Twitter for Python! \\
\end{longtable}
\end{CJK*}
% \captionof{table}{List of all GitHub repositories represented in \bugs{} and their corresponding descriptions. We identify 7 categories generally characterizing the range of utilities for the selected codebases.}
\label{tab:summary_of_repos}

\subsection{Bug Generation Statistics}
\label{appx:dataset:bug_gen_stats}
We provide extensive details about different aspects of each of the bug generation strategies, including the yield rates, labor/monetary costs, and dataset characterizations.

\textbf{Yield rates.}
In Table~\ref{tab:per_bug_type_yield_rate}, we provide the yield rates for each bug generation method across all repositories in \bugs{}.
In general, we find that the PR Mirroring has the lowest yield rate at $13.18$\% (although this rate is somewhat higher than SWE-bench's yield rate of $2294/93139 = 2.46$\%).
For using LMs to generate bugs, modifying functions to introduce bugs intentionally has a higher yield than asking LMs to perform a best-effort rewrite.
The efficacy of Procedural Modifications varies by strategy.
For instance, shuffling the functions declared in a class only breaks existing test(s) $1.93$\% of the time, but inverting a conditional will lead to a task instance for $47.04$\% of modifications.
Finally, combining bug patches has an extremely high yield rate - this is to be expected because we only attempt to combine bug patches that have been validated as usable task instances breaking $1+$ tests.

\begin{table}[h]
    \centering
    \begin{tabular}{l|rrrr}
\toprule
Strategy & \# Repos & \# Candidates & \# Instances & Yield Rate \\
\midrule
Combine (file) & 124 & 6020 & 5865 & 97.43\% \\
Combine (module) & 65 & 4396 & 4227 & 96.16\% \\
LM (Modify) & 108 & 31950 & 17887 & 55.98\% \\
LM (Rewrite) & 128 & 11908 & 4173 & 35.04\% \\
PR Mirroring & 108 & 6934 & 2344 & 33.8\% \\
Procedural (Class Rm Base) & 103 & 1401 & 463 & 33.05\% \\
Procedural (Class Rm Funcs) & 103 & 2506 & 1180 & 47.09\% \\
Procedural (Class Shuffle Funcs) & 103 & 2504 & 47 & 1.88\% \\
Procedural (Ctrl Invert If) & 105 & 4695 & 2321 & 49.44\% \\
Procedural (Ctrl Shuffle) & 104 & 9055 & 4015 & 44.34\% \\
Procedural (Op Break Chains) & 71 & 747 & 225 & 30.12\% \\
Procedural (Op Change Const) & 77 & 723 & 257 & 35.55\% \\
Procedural (Op Change) & 81 & 1507 & 450 & 29.86\% \\
Procedural (Op Swap) & 87 & 2141 & 483 & 22.56\% \\
Procedural (Remove Assign) & 121 & 5470 & 2661 & 48.65\% \\
Procedural (Remove Cond) & 120 & 5288 & 2311 & 43.7\% \\
Procedural (Remove Loop) & 110 & 1945 & 860 & 44.22\% \\
Procedural (Remove Wrapper) & 80 & 884 & 368 & 41.63\% \\
\midrule
All & 129 & 100074 & 50137 & 50.1\% \\
\bottomrule
    \end{tabular}
    \caption{
    Yield rates for different bug generation strategies covered in Section~\ref{appx:bugs:generate}.
    We show the number of repositories that each strategy was run on, the number of bug candidates generated by each strategy, and the number of instances, or the number of candidates that were validated to have $1+$ Fail to Pass test.
    The yield rate for 
    }
    \label{tab:per_bug_type_yield_rate}
\end{table}

The number of repositories captured by each bug generation technique varies due to each strategy's specific preconditions, which at times may not be effective for some repositories.
For instance, the \textit{Procedural (Class *)} set of methods only mutates Python classes.
This strategy is fruitless for the minority of \bugs{} repositories that do not define any classes.
The \textit{Procedural (Op Break Chains)} method randomly removes operations and operands from expressions with two or more operations (e.g. $a+b+c \rightarrow a+b$) --- such expressions are not always present in \bugs{} repositories.

The collective yield rate across \bugs{}'s bug generation strategies is significantly higher than SWE-bench's collection strategy.

\begin{table}[ht]
    \centering
    \begin{tabular}{l|c}
\toprule
Yield Rate & \# of Repositories \\
\midrule
$0$-$25$\%   & $10$  \\
$25$-$50$\%  & $31$ \\
$50$-$75$\%  & $60$ \\
$75$-$100$\% & $27$ \\
\bottomrule
    \end{tabular}
    \caption{
    Yield rates for different repositories represented in \bugs{}.
    }
    \label{tab:per_repo_yield_rate}
\end{table}

The yield rate also varies with respect to the repository it is being applied to.
We provide a summary of yield rates by repository in Table~\ref{tab:per_repo_yield_rate}.
We generally observe that lower test coverage correlates with a lower yield rate.

\textbf{Dataset characterizations.}
In Table~\ref{tab:per_bug_type_stats}, we provide statistics about the validated task instances produced by different bug generation strategies.
Our work's LM-based strategies rewrite one function in one file.
Procedural modifications will also only change one file, but depending on the strategy, $1+$ functions or classes may be changed.
Combining multiple patches from the same file always produces a patch with $2+$ functions edited.
Combining across modules produces a patch with $2+$ files edited.
The targeted nature of each of the bug creation strategies is reflected in the typical number of functions and files that the bugs produced by each strategy edits.

\begin{table}[h]
    \centering
    \begin{tabular}{l|rrrrr}
\toprule
Strategy & \# Instances & \# F2P & $\Delta$ Lines & $\Delta$ Functions & $\Delta$ Files \\
\midrule
Combine      & 10092 & 15 (5-48) & 19 (12-36) & 2 (2-3) & 1 (1-2) \\
LM           & 22060 & 4 (1-17) & 6 (3-15) & 1 (1-1) & 1 (1-1) \\
PR Mirroring & 2344 & 3 (1-14) & 20 (8-55) & 2 (2-4) & 1 (1-2) \\
Procedural   & 15641 & 7 (2-32) & 7 (5-15) & 1 (1-1) & 1 (1-1) \\
\bottomrule
    \end{tabular}
    \caption{Statistics for attributes of a \bugs{} task instance across different bug generation strategies, reported as \textit{median (IQR)}, where IQR is the inter-quartile range (25th–75th percentile).
    }
    \label{tab:per_bug_type_stats}
\end{table}

In Figure~\ref{fig:dataset_comparison}, we show the distributions for different attributes of \bugs{} compared to other SWE-bench style datasets.
Compared to prior works, there is a much higher proportion of task instances with more than one Fail-to-Pass test.
For any one repository, we find that \bugs{} task instances collectively cause failures for a much higher percentage of the testing suit than other datasets; a potential benefit of this is that training on \bugs{} based trajectories may expose models to a much broader set of functionalities in a codebase.
The number of lines and files edited by \bugs{} task instances is highly similar to the trend lines for SWE-bench Verified.

\begin{figure}[h]
    \centering
    \includegraphics[width=\textwidth]{figures/dataset_comparison.pdf}
    \caption{
    Comparison of cumulative distributions for Fail-to-Pass tests along with the lines and files edited by the gold patch across \bugs{} and four SWE-bench style datasets.
    }
    \label{fig:dataset_comparison}
\end{figure}

We note that unlike other datasets, the trend line of \bugs{} task instances is ``adjustable".
In other words, the Figure~\ref{fig:dataset_comparison} distributions are a capture of the task instances provided in this release of \bugs{}.
However, because of \bugs{}'s flexible bug creation techniques, the distribution can be ``shaped" if needed.
For instance, generating more task instances using the bug patch combination method would shift all three curves in Figure~\ref{fig:dataset_comparison}.
We make this point to highlight the fact that the attributes of SWE-bench task instances are, in a sense, constrained by real world software development behavior.
On the other hand, \bugs{} can be used to break tests and code that may not be reflected at all in any existing pull request.
In this sense, we argue that LMs trained on \bugs{} have better ``exposure" to a codebase compared to exclusively training on pull requests.

\textbf{Continuation of scaling execution environments.}
The validation and evaluation procedures for \bugs{} deviate slightly from SWE-bench's harnesses.
The main reasons for these differences can largely be attributed to the granularity of installation specifications.
In SWE-bench, each task instance corresponds to a unique base commit, with additional \texttt{version} and \texttt{environment\_setup\_commit} keys needed as indirection for mapping an instance to the correct set of installation and testing instructions.
Across time, the continuous evolution of a repository and its dependencies make for an incredibly high degree of variability in how a repository should be installed correctly.
To solve this variability, the community has resorted to creating an image per task instance, as done in~\citet{chowdhuryintroducing}.
Therefore, for $2294$ SWE-bench task instances, there are $2294$ unique Docker images, each at a size of at least several gigabytes ($\sim5$-$6$ GBs).

On the other hand, the simplicity and scalability of \bugs{}'s design allows one to support many task instances with comparatively much fewer Docker images.
As mentioned above, installation and testing procedures are (repository, commit) specific.
Therefore, when bugs are generated from each (repository, commit), all bugs can be reproduced and tested successfully from the same Docker image.
In other words, if I generate $100$ bugs for a repository at some commit, instead of $100$ Docker images, only a single Docker image is required to run inference on any of the $100$ task instances.

This design is what enables \bugs{} to be significantly more space-efficient than SWE-bench.
Based on the publicly released images, for SWE-bench's $2294$ task instances, $1.2$ TBs of storage are required to download all Docker images locally.
for SWE-bench Multimodal's $517$ task instances, $1.2$ TBs are required.
The higher per-instance Docker image size for SWE-bench Multimodal is due to how JavaScript dependency management tools (e.g. \texttt{npm}) require more storage compared to equivalent Python infrastructure (e.g. \texttt{pypi}).
\citet{pan_training_2024} states that each image for the $2438$ instances an average of $2.6$GB, totaling 6 TB of storage total.
Such a storage requirement can be a significant barrier for academic practitioners.

On the other hand, with more than $20$x the number of bugs, \bugs{} requires only $125$ Docker images total, corresponding to the number of unique (repository, commit) pairs (in this work, for each repository, we only determine installation and test specifications for one commit).
The $125$ images require a total of $290.54$ GBs.
In summary, compared to SWE-bench's task collection strategy, \bugs{}'s design makes it easier to not only create task instances, but also train on them as well.

\subsection{Case Study: SWE-bench \& \bugs{}}
\label{appx:dataset:case-study}
To better understand the differences between the SWE-bench and \bugs{} collection strategies, we perform \bugs{} collection on the \texttt{pallets/flask} GitHub repository, one of the $12$ test split repositories from the original SWE-bench benchmark.
We review the steps covered in Section~\ref{sec:swesmith:collection} applied to \texttt{pallets/flask} in detail.
First, we defined the installation and testing specifications for the \texttt{pallets/flask} repository at commit \texttt{bc09840}.
Next, we apply the LM modification bug generation strategy to this version of the repository, generating $267$ unique bugs.

We observe several differences.
First, \textit{the \bugs{} collection strategy yields a much higher number of bugs outright.}
From SWE-bench, $11$ task instances are from the \texttt{pallets/flask} repository.
The task instances were originally filtered from $2434$ pull requests (PRs), with $107$ satisfying SWE-bench's filtering criteria of (1) being linked to one or more issues and (2) featuring 1+ new tests.
Out of these $107$, the $11$ ($0.45$\% of $2434$) task instances represent the proportion of PRs that execution environments could be successfully constructed for.
On the other hand, running the function-level rewriting strategy for bug generation originally yielded $402$ candidates, of which $267$ were determined to be valid task instances.

Second, \textit{\bugs{} requires significantly less human effort while only incurring minor costs}.
Collecting the $11$ \texttt{pallets/flask} task instances (steps include scraping PRs, determining repository versions across time, defining version-specific installation/test specifications, running execution-based validation multiple times) took an estimated $38$ hours worth of human labor.
On the contrary, defining installation and testing specifications for the latest commit of \texttt{pallets/flasks} took $10$ minutes. 
The subsequent function-level rewriting strategy for bugs took $23$ minutes to run, incurring a total cost of just \$$2.47$ ($\sim$\$$0.00613$ per instance).
The final execution-based validation step that filters out $402 - 267 = 135$ unqualified bug candidates ran in $14$ minutes.
Since both the bug and problem statement generation strategies are repository agnostic, no additional human intervention is necessary for these steps.
Head to head, per instance for the \texttt{pallets/flask} repository, SWE-bench style collection requires $38 \times 60 / 11 = 207.27$ minutes compared to $0.176$ minutes ($\sim10.6$ seconds) and \$$0.00613$ in API costs using \bugs{}.

Third, \textit{collectively, \bugs{} task instances break a significantly larger proportion of existing tests in a codebase}.
We define ``bug coverage" as the proportion of tests broken by $1$+ instance across all task instances.
For the SWE-bench split of \texttt{pallets/flask}, there are $207$ unique tests across all $11$ instances.
Of these $207$ tests, $15$ are broken by $1$+ instance, corresponding to a bug coverage rate of $7.25$\%.
For the \bugs{} split of \texttt{pallets/flask}, there are $474$ unique tests across $267$ instances.
The larger amount of tests is due to increased test coverage in the \texttt{pallets/flask} repository as of Nov. 28, 2024 (when \bugs{} was collected) compared to June 2023 (when SWE-bench was collected).
Of these $474$ tests, $422$ are broken by $1$+ instance, a bug coverage rate of $89.03$\%.
We attribute the significant difference to a consistent tendency in real world open source software development workflows, that is, the \textit{minority} of tests are introduced to capture existing, errant behavior in the repository.
The significant majority of tests are committed alongside working code, ensuring that already correct behavior is upheld.
Well-maintained repositories will typically not merge commits that cause such tests to fail.
This results in a large number of tests where few to no commits correspond to those tests' failures.

Finally, \textit{\bugs{} does not yield instances appropriate for evaluation}.
The \bugs{} pipeline as presented does not produce hidden tests, a crucial difference that makes SWE-bench more suitable for evaluation.
Consequently, when expert trajectories are generated, the Fail-to-Pass tests are present in the repository at inference time.
Furthermore, our issue generation strategy does not include checks for known problems such as underspecified text descriptions or solution leakage~\citep{chowdhuryintroducing}. 
Simple amendments could make \bugs{} task instances suitable for evaluation, such as deleting Fail-to-Pass test functions or files along with a validation procedure around the ambiguity and leakage of the issue text.
Finally, thorough analyses of how faithful \bugs{} task instances are to real world issues and PRs would be necessary to justify synthetic bugs for evaluation.
